% Môn: Vật lý
% Lớp: 11
% Chương: Chương I. Dao động
% Bài: L11C1 Bài 6. Dao động tắt dần. Dao động cưỡng bức. Hiện tượng cộng hưởng
% Số câu: 9
% App đọc file này trực tiếp — sửa rồi Commit trên GitHub, bấm Đồng bộ GitHub .tex

% L11C1 Bài 6. Dao động tắt dần. Dao động cưỡng bức. Hiện tượng cộng hưởng

% ===== Câu 1 =====
% ID: L11C1B6-01-TL
% Mức: NB
\dangbt{Chưa có dạng}
\begin{ex}
% ID: L11C1B6-01-TL
% Mức: NB
Một con lắc lò xo gồm lò xo có khối lượng không đáng kể. Chu kì dao động của con lắc là 0,1 $\pi(s).$ Con lắc dao động cưỡng bức theo phương trùng với trục của lò xo dưới tác dụng của ngoại lực tuần hoàn $F = F_0 \cos \omega t(𝑁)$. Khi $\omega$ lần lượt là 10rad/s và 15rad/s thì biên độ dao động tương ứng của con lắc lần lượt là A_1 và A_2. Hãy so sánh A_1 và A_2. $\nguon{ly11 kntt.pdf}$
\choiceTF

\loigiai{
Khi xảy ra hiện tượng cộng hưởng: $\omega_{0} = \dfrac{2\pi}{T} = \dfrac{2\pi}{0,1\pi} = 20\text{rad}/\text{s}$ Vì ω₁ = $10\,\mathrm{rad/s}$ xa vị trí cộng hưởng hơn ω₂ = $15\,\mathrm{rad/s}$ (ω₁<ω₂<ω₀) nên A₁ < A₂ (Hình $6.1\,\text{G}$).
}
\end{ex}

% ===== Câu 2 =====
% ID: L11C1B6-01-TN
% Mức: NB
\begin{ex}
% ID: L11C1B6-01-TN
% Mức: NB
Trong dao động tắt dần, một phần cơ năng đã biến đổi thành $\nguon{ly11 kntt.pdf}$
\choice
{điện năng.}
{\True nhiệt năng.}
{hoá năng.}
{quang năng.}
\loigiai{
Chọn B.

Trong dao động tắt dần, một phần cơ năng đã biến đổi thành nhiệt năng.
}
\end{ex}

% ===== Câu 3 =====
% ID: L11C1B6-02-TL
% Mức: TH
\begin{bt}
% ID: L11C1B6-02-TL
% Mức: TH
Một con lắc lò xo gồm vật nhỏ khối lượng $m = 0$, $2 \mathrm{kg}$, lò xo nhẹ có độ cứng $k = 20\$, $\mathrm{N}/m$ dao động trên mặt phẳng nằm ngang. Hệ số ma sát trượt giữa vật và mặt phẳng nằm ngang là $\mu = 0$,01. Từ vị trí lò xo không bị biến dạng, truyền cho vật vận tốc ban đầu có độ lớn $v_0 =$ 1 $\mathrm{m/s}$ dọc theo trục lò xo (lấy g = 10 m/s_2 ) $. Tính độ lớn lực đàn hồi cực đại của lò xo trong quá trình dao động.$ \nguon{ly11kntt.pdf}
\loigiai{
Độ lớn của lực đàn hồi sẽ đạt cực đại khi vật ra tới vị trí biên lần đầu tiên sau khi được truyền vận tốc v₀ (vì biên độ ở các lần sau sẽ không bằng được ở lần này). Công của lực ma sát trên đoạn biên độ $A$ đầu tiên đó bằng độ giảm cơ năng khi vật đi từ vị trí ban đầu tới vị trí biên: $- \mu\text{mgA} = \dfrac{kA^{2}}{2} - \dfrac{mv_{0}^{2}}{2}$

Thay số: $. - 0,01 \cdot 0,2 \cdot$ 10 $\mathrm{A}$ = \dfrac{ $20$ \mathrm{A} $^{2}}{2} - \dfrac{0,2 \cdot 1^{2}}{2}\Rightarrow$ 10 $\mathrm{A}$ ^{2} + 0, $02$ \mathrm{A} $- 0,1 = 0 .\Rightarrow$ A=0, $099\,\mathrm{m}\RightarrowF_(dhmax)=kA=20.0,099=1,$ 98\,\mathrm{N}.
}
\end{bt}

% ===== Câu 4 =====
% ID: L11C1B6-02-TN
% Mức: NB
\begin{ex}
% ID: L11C1B6-02-TN
% Mức: NB
Một con lắc lò xo đang dao động tắt dần, sau ba chu kì đầu tiên, biên độ của nó giảm đi 10%. Phần trăm cơ năng còn lại sau khoảng thời gian đó là $\nguon{ly11 kntt.pdf}$
\choice
{\True 81%.}
{6,3%.}
{19%.}
{27%.}
\loigiai{
Chọn A.

Theo đề bài: $\left. \dfrac{A - A_{3}}{A} = 10\text{\%}\Rightarrow\dfrac{A_{3}}{A} = 90\text{\%}\Rightarrow\dfrac{W_{3}}{W} = \left( \dfrac{A_{3}}{A} \right)^{2} = 0,9^{2} = 0,81 = 81\text{\%} \right.$
}
\end{ex}

% ===== Câu 5 =====
% ID: L11C1B6-03-TL
% Mức: TH
\begin{ex}
% ID: L11C1B6-03-TL
% Mức: TH
Một con lắc lò xo gồm vật nhỏ có khối lượng $m = 0$, $03 \mathrm{kg}$ và lò xo có độ cứng $k = 1$, $5\,\mathrm{N}$ /m. Vật nhỏ được đặt trên giá đỡ cố định nằm ngang dọc theo trục của lò xo. Hệ số ma sát trượt giữa giá đỡ và vật nhỏ là $\mu = 0$,2. Ban đầu, giữ vật ở vị trí lò xo bị dãn một đoạn $\Deltal_0 =$ 15 $\mathrm{cm}$ rồi buông nhẹ để con lắc dao động tắt dần $. Lấy$ g = 10 m/s_2 $. Tính tốc độ lớn nhất mà vật nhỏ đạt được trong quá trình dao động.$ \nguon{ly11kntt.pdf}
\choiceTF

\loigiai{
Vật đạt tốc độ lớn nhất tại vị trí $O$ mà lực ma sát cân bằng với lực đàn hồi của lò xo, khi đó vật còn cách vị trí mà lò xo không biến dạng một đoạn $\Delta$ l xác định bởi: $\left. \mu\text{mg} = \text{k} \cdot \text{\Delta }l\Rightarrow\text{\Delta }l = \dfrac{\mu\text{mg}}{\text{k}} = \dfrac{0,2 \cdot 0,03 \cdot 10}{1,5} = 0,04\text{m.} \right.$ [Một con lắc lò xo gồm vật nhỏ có khối lượng m = $0,03\,\text{kg}$]

Công của lực ma sát trên đoạn $\Delta$ l₀ - $\Delta$ l đó bằng độ giảm cơ năng khi vật đi từ vị trí ban đầu tới vị trí cân bằng nói trên:

$- \mu\text{mg}\left( {\text{\Delta }l_{0} - \text{\Delta }l} \right) = \dfrac{\text{mv}_{\max}^{2}}{2} + \dfrac{\text{k} \cdot \text{\Delta }l^{2}}{2} - \dfrac{\text{k} \cdot \text{\Delta }l_{0}^{2}}{2}$ Thay số: -0,1.0,03.10(0,15-0,04) = $\dfrac{0,03v_{\text{max}}^{2}}{2} + \dfrac{1,5 \cdot 0,04^{2}}{2} - \dfrac{1,5 \cdot 0,15^{2}}{2}$ Suy ra: v_(max) = 0, $91\,\mathrm{m}$ /s = $91\,\mathrm{cm}$ /s.
}
\end{ex}

% ===== Câu 6 =====
% ID: L11C1B6-03-TN
% Mức: NB
\begin{ex}
% ID: L11C1B6-03-TN
% Mức: NB
Một con lắc lò xo dao động tắt dần theo phương ngang với chu kì $T = 0$, $2\,\mathrm{s}$, lò xo nhẹ, vật nhỏ dao động có khối lượng $100\,\mathrm{g}$. Hệ số ma sát giữa vật và mặt phẳng ngang là 0,01. Độ giảm biên độ của vật sau mỗi lần vật đi từ biên này tới biên kia là $\nguon{ly11 kntt.pdf}$
\choice
{$0,02 \mathrm{mm}$.}
{$0,04 \mathrm{mm}$.}
{\True $0,2 \mathrm{mm}$.}
{$0,4 \mathrm{mm}$.}
\loigiai{
Chọn C.

Nguyên nhân của dao động tắt dần là do lực cản của môi trường, trong bài toán này là lực ma sát. Độ giảm cơ năng sau một nửa chu kì bằng công của lực ma sát thực hiện trong chu kì đó, ta có:

$\left. \dfrac{1}{2}m\omega^{2}A^{2} - \dfrac{1}{2}m\omega^{2}A^{\text{'}2} = F_{ms}\left( {A + A^{'}} \right)\Leftrightarrow\dfrac{1}{2}m\omega^{2}\left( {A + A^{'}} \right)\left( {A - A^{'}} \right) = F_{ms}\left( {A + A^{'}} \right) \right.\left. \Rightarrow\text{\Delta }A = \dfrac{2F_{ms}}{k} = \dfrac{2\mu mg}{k} \right.$

Độ giảm biên độ sau mỗi lần qua vị trí cân bằng:

$\dfrac{\text{\Delta A}}{2} = \dfrac{2\mu\text{mg}}{\text{k}} = \dfrac{2 \cdot 0,01 \cdot 0,1 \cdot 10}{100} = 0,2 \cdot 10^{-3}\text{m}.$
}
\end{ex}

% ===== Câu 7 =====
% ID: L11C1B6-04-TL
% Mức: TH
\begin{ex}
% ID: L11C1B6-04-TL
% Mức: TH
Một con lắc lò xo gồm vật nhỏ khối lượng $m = 0$, $02 \mathrm{kg}$ và lò xo có độ cứng $k = 1\$, $\mathrm{N}/m.$ Vật nhỏ được đặt trên giá đỡ cố định nằm ngang dọc theo trục lò xo. Hệ số ma sát trượt giữa giá đỡ và vật nhỏ là $\mu = 0$,1. Ban đầu giữ vật ở vị trí lò xo bị nén $\Deltal_0 =$ 10 $\mathrm{cm}$ rồi buông nhẹ để con lắc dao động tắt dần $. Lấy$ g = $10$ \mathrm{m/s} $2$. Tính độ giảm thế năng của con lắc trong giai đoạn từ khi buông tới vị trí mà tốc độ dao động của con lắc cực đại lần đầu. BÀ $I$ TẬ $P$ CUỐ $I$ CHƯƠNG $I\nguon{ly11 kntt.pdf}$
\choiceTF

\loigiai{
Vật đạt tốc độ lớn nhất v_(max) tại vị trí mà F_(ms) = F_(oh) (Hình $6.2\,\text{G}$)

[Một con lắc lò xo gồm vật nhỏ có khối lượng m = $0,02\,\text{kg}$]

⇔μmg=k $\Delta$ l

$\left. \Rightarrow\text{\Delta }l = \dfrac{\mu\text{mg}}{\text{k}} = \dfrac{0,1 \cdot 0,02 \cdot 10}{1} = 0,02\text{m} \right.\left. \Rightarrow\text{\Delta W}_{\text{t}} = \dfrac{\text{k} \cdot \text{\Delta }l_{0}^{2}}{2} - \dfrac{\text{k} \cdot \text{\Delta }l^{2}}{2} = \dfrac{\text{k}}{2}\left( {\text{\Delta }l_{0}^{2} - \text{\Delta }l^{2}} \right) \right.= \dfrac{1}{2}\left( {0,1^{2} - 0,02^{2}} \right) = 4,8 \cdot 10^{-3}\text{J} = 4,8\text{mJ}$
}
\end{ex}

% ===== Câu 8 =====
% ID: L11C1B6-04-TN
% Mức: NB
\begin{ex}
% ID: L11C1B6-04-TN
% Mức: NB
Một người xách một xô nước đi trên đường, mỗi bước đi dài $L =$ 50 $\mathrm{cm}$ thì nước $trong xô bị sóng sánh mạnh nhất. Tốc độ đi của người đó là$ v = 2 $,$ 5 \mathrm{km} $/h. Chu kì dao động riêng của nước trong xô là$ \nguon{ly11kntt.pdf}
\choice
{\True $1,44\,\mathrm{s}$.}
{$0,35\,\mathrm{s}$.}
{$0,45\,\mathrm{s}$.}
{$0,52\,\mathrm{s}$.}
\loigiai{
Chọn A.

Hai bước đi là một chu kì. Chu kì dao động riêng của nước trong xô là:

$\text{T} = \dfrac{2L}{\text{V}} = \dfrac{2 \cdot 0,5}{0,69} \approx 1,44\text{s.}$
}
\end{ex}

% ===== Câu 9 =====
% ID: L11C1B6-05-TN
% Mức: TH
\begin{ex}
% ID: L11C1B6-05-TN
% Mức: TH
Tìm phát biểu sai. Dao động tắt dần là dao động có $\nguon{ly11 kntt.pdf}$
\choice
{\True tần số giảm dần theo thời gian.}
{cơ năng giảm dần theo thời gian.}
{biên độ dao động giảm dần theo thời gian.}
{ma sát và lực cản càng lớn thì dao động tắt dần càng nhanh.}
\loigiai{
Chọn A.

Dao động tắt dần là dao động có biên độ, cơ năng giảm dần theo thời gian, khi ma sát và lực cản càng lớn thì dao động tắt dần càng nhanh.
}
\end{ex}
